\section{o1 和 R1 的真正启示}

\subsection{RL 扩展的关键前提}

第一波 reasoning 模型告诉我们一个根本性的道理：如果我们想在语言模型中扩展强化学习，就需要\textbf{确定性的、稳定的、可扩展的反馈信号}。数学、代码、逻辑等可验证领域因此成为核心训练场景，因为这些领域的奖励信号比通用的偏好监督强得多。它们让 RL 可以针对\textbf{正确性}（correctness）而非\textbf{似然性}（plausibility）进行优化。

\begin{knowledgebox}{从 Pretraining Scaling 到 Post-training Scaling}
o1 和 R1 开启了一个关键转型：从扩展预训练（scaling pretraining）到扩展后训练中的 reasoning 能力（scaling post-training for reasoning）。这是 LLM 发展路径上的第一次大转折。
\end{knowledgebox}

\subsection{基础设施成为核心问题}

一旦模型被训练在更长的轨迹上进行推理，RL 就不再是 SFT 之上的轻量附加组件，而变成了一个\textbf{系统工程问题}。这需要大规模的 rollout、高吞吐量的验证、稳定的策略更新和高效的采样。正如作者所言：\textit{The emergence of reasoning models was as much an infra story as a modeling story.}（reasoning 模型的兴起既是基础设施的故事，也是建模的故事。）

\subsection{本章小结}

o1/R1 时代的核心教训有两点：（1）RL 扩展需要强反馈信号，可验证领域是理想训练场；（2）reasoning RL 本质上是一个系统工程问题，基础设施的重要性与模型本身相当。


